\documentclass[12pt,a4paper,oneside]{article}
\usepackage{fontspec}
\setmainfont{Times New Roman}
\setmonofont{Courier New}
\usepackage[english,russian]{babel}
\usepackage{geometry}
\geometry{left=2.54cm,right=2.54cm,top=2.54cm,bottom=2.54cm}
\usepackage{amsmath,amssymb,array,tabularx,booktabs}
\usepackage{graphicx,listings,xcolor,hyperref,enumitem,titlesec}
\titleformat*{\section}{\large\bfseries}
\titleformat*{\subsection}{\normalsize\bfseries}
\setcounter{secnumdepth}{0}
\hypersetup{pdftitle={Отчёт по ЛР № 4. ОС. Вариант 18},pdfauthor={}}
\setlength{\parindent}{0pt}
\setlength{\parskip}{5pt}
\sloppy
\lstset{basicstyle=\footnotesize\ttfamily,frame=single,breaklines=true,
  columns=fullflexible,numbers=left,numberstyle=\tiny,
  showstringspaces=false,tabsize=4,keepspaces=true}
\newcommand{\code}[1]{\texttt{\detokenize{#1}}}
\newcommand{\newsection}[1]{\section{#1}}
\newcommand{\smalltable}{\small}
\begin{document}
\begin{titlepage}
\begin{center}
\large МОСКОВСКИЙ АВИАЦИОННЫЙ ИНСТИТУТ\\
(НАЦИОНАЛЬНЫЙ ИССЛЕДОВАТЕЛЬСКИЙ УНИВЕРСИТЕТ)
\vspace{20pt}

Институт №8 «Компьютерные науки и прикладная математика»\\
Кафедра 806 «Вычислительная математика и программирование»
\end{center}
\vspace{65pt}
\begin{center}
\large\bfseries Отчет по лабораторной работе №4\\[7pt]
по курсу «Операционные системы»
\end{center}
\vspace{40pt}
\begin{center}
\large\bfseries Тема: «Межпроцессное взаимодействие через отображаемую память»\\[7pt]
Вариант 18
\end{center}
\vfill
\begin{flushright}
\begin{tabular}{rl}
Выполнил(а): & А.\,Е. Пронина \\
Группа: & М8О-203БВ-25 \\
Преподаватель: & Е.\,С. Миронов \\
\end{tabular}
\end{flushright}
\vspace{40pt}
\begin{center}Москва, 2026\end{center}
\end{titlepage}

\section{Цель работы}
Изучить работу с файловыми объектами и организацию обмена данными между процессами с помощью отображаемой памяти (file mapping). Реализовать передачу строк через \code{mmap()}, синхронизировать доступ POSIX-семафорами и обеспечить освобождение ресурсов при штатном завершении.

\section{Постановка задачи}
Для варианта 18 требуется реализовать родительский процесс и два дочерних процесса. Родитель считывает имена двух выходных файлов и строки стандартного ввода. Нечётные строки данных направляются первому ребёнку, чётные --- второму. Пустая строка также имеет номер и участвует в чередовании.

Каждый дочерний процесс удаляет из полученных строк английские гласные и записывает результат в свой файл. В исследуемом коде удаляются символы \code{a, e, i, o, u} и их прописные варианты. Символ \code{y} и русские гласные данным фильтром не удаляются.

Для ЛР № 4 обмен организуется через две области общей памяти. Дочерние процессы создаются с помощью \code{fork()}, после чего выполняется \code{execvp()} программы \code{child}. После EOF родитель передаёт признак завершения, ожидает детей и освобождает объекты общей памяти.

\section{Архитектура решения}
Программа состоит из двух исполняемых файлов: \code{parent} и \code{child}. Каждый ребёнок обслуживает собственный канал с единственным буфером.

\begin{center}
\fbox{\parbox{0.88\linewidth}{\centering
Стандартный ввод $\longrightarrow$ \code{parent}\\[3pt]
нечётные строки $\longrightarrow$ память 1 $\longrightarrow$ \code{child 1} $\longrightarrow$ файл 1\\[3pt]
чётные строки $\longrightarrow$ память 2 $\longrightarrow$ \code{child 2} $\longrightarrow$ файл 2\\[3pt]
В каждой области: данные, флаги и два POSIX-семафора
}}
\end{center}
Служебные сообщения дочерних процессов выводятся в \code{stderr}. Результат фильтрации направляется в \code{stdout}, который ребёнок перенаправляет в выходной файл через \code{open()} и \code{dup2()}.

\newsection{Общая память и протокол обмена}
\subsection{Структура данных}
В \code{shared.h} объявлена следующая структура:
\begin{lstlisting}[language=C]
#define BUFFER_SIZE 4096

typedef struct {
    sem_t can_write;
    sem_t can_read;
    size_t length;
    int end_of_line;
    int finished;
    char text[BUFFER_SIZE];
} SharedData;
\end{lstlisting}
\code{length} задаёт число передаваемых байтов; \code{end_of_line} отмечает последнюю часть строки; \code{finished} сообщает об окончании обмена. Буфер не обязан содержать завершающий нулевой байт: фильтр и \code{fwrite()} используют явную длину.

Объекты именуются по шаблону \code{/lab4_<PID>_1} и \code{/lab4_<PID>_2}. В трассировке размер каждого отображения равен 4176 байтам: это размер всей структуры в данной сборке, а не только массива \code{text}. Размер полезной порции равен 4096 байтам.

\subsection{Синхронизация}
Оба семафора размещаются непосредственно в общей памяти. Параметр \code{pshared=1} при \code{sem_init()} разрешает межпроцессное использование. Начальные значения: \code{can_write=1}, \code{can_read=0}.
\begin{enumerate}
\item Родитель ожидает \code{can_write}, записывает данные, длину и флаги, затем выполняет \code{sem_post(can_read)}.
\item Ребёнок ожидает \code{can_read}, обрабатывает порцию и записывает результат. После этого он выполняет \code{sem_post(can_write)}.
\item Следующая порция перезаписывает буфер только после разрешения ребёнка. Между частями одной строки перевод строки не добавляется.
\end{enumerate}
При EOF отправляется сообщение с \code{finished=1}. Получив его, ребёнок освобождает семафор записи и завершает цикл.

\newsection{Используемые механизмы ОС}
\subsection{Отображение памяти}
\code{shm_open()} создаёт или открывает именованный объект POSIX shared memory. \code{ftruncate()} задаёт его размер, а \code{mmap()} отображает объект в виртуальное адресное пространство процесса. После успешного отображения файловый дескриптор закрывается; само отображение продолжает существовать.

Процессы имеют отдельные виртуальные адресные пространства. Поэтому один объект может находиться по разным виртуальным адресам у родителя и ребёнка. Флаг \code{MAP_SHARED} обеспечивает совместный доступ к изменениям объекта. При \code{MAP_PRIVATE} изменения через отображение не использовались бы как общий канал записи между процессами.

Одного \code{MAP_SHARED} недостаточно для согласования доступа: семафоры задают порядок заполнения и чтения буфера. Они предотвращают чтение частично подготовленного сообщения и перезапись ещё не обработанных данных.

\subsection{Распределение и обработка строк}
В \code{parent.c} строка читается функцией \code{getline()}, завершающий \code{\n} исключается из полезных данных. Выбор ребёнка выполняется по остатку от деления счётчика \code{line_number} на 2. Счётчик начинается с нуля, поэтому первая строка направляется ребёнку 1.

\code{send_line()} разбивает строку на части до 4096 байт. Цикл \code{do ... while} отправляет сообщение даже для пустой строки: длина равна нулю, а \code{end_of_line=1}. Ребёнок в этом случае записывает один перевод строки.

\code{remove_vowels()} проходит по данным один раз и уплотняет сохраняемые символы в том же буфере. Для строки длины $n$ обработка занимает $O(n)$ времени и $O(1)$ дополнительной памяти фильтра. Память родителя для \code{getline()} зависит от длины строки.

\newsection{Ключевые фрагменты исходного кода}
Ниже приведены ключевые фрагменты программы. Обозначение \code{/* ... */} показывает сокращённые участки обработки ошибок.

\subsection{Создание области общей памяти}
Основные операции функции \code{create_shared()} из \code{shared.c}:
\begin{lstlisting}[language=C]
int fd = shm_open(name, O_CREAT | O_EXCL | O_RDWR, 0600);
/* ... */
if (ftruncate(fd, sizeof(SharedData)) == -1) {
    perror("ftruncate");
    cleanup_shared(name, fd, MAP_FAILED, 0, 0);
    return -1;
}
void* ptr = mmap(NULL, sizeof(SharedData),
    PROT_READ | PROT_WRITE, MAP_SHARED, fd, 0);
/* ... */
SharedData* shared = (SharedData*)ptr;
if (sem_init(&shared->can_write, 1, 1) == -1) {
    /* ... */
}
if (sem_init(&shared->can_read, 1, 0) == -1) {
    /* ... */
}
\end{lstlisting}
\code{O_EXCL} предотвращает открытие уже существующего объекта под видом нового. Функция \code{cleanup_shared()} освобождает успешно созданные части ресурса при ошибке и удаляет имя объекта.

\subsection{Публикация сообщения}
В \code{send_message()} родитель ожидает \code{can_write} через \code{sem_timedwait()} с очередным сроком ожидания через одну секунду. После \code{EINTR} или \code{ETIMEDOUT} выполняется \code{waitpid(..., WNOHANG)}: обнаруженное завершение ребёнка прерывает передачу. После успешного ожидания выполняется:
\begin{lstlisting}[language=C]
if (length > 0) memcpy(child->data->text, text, length);
child->data->length = length;
child->data->end_of_line = end_of_line;
child->data->finished = finished;
if (sem_post(&child->data->can_read) == -1) {
    perror("sem_post");
    return -1;
}
\end{lstlisting}

\subsection{Передача строки частями}
Функция \code{send_line()} из \code{parent.c} приведена полностью:
\begin{lstlisting}[language=C]
static int send_line(ChildInfo* child,
                     const char* line, size_t length) {
    size_t offset = 0;
    do {
        size_t part = length - offset;
        if (part > BUFFER_SIZE) part = BUFFER_SIZE;
        int last = (offset + part == length);
        if (send_message(child, line + offset,
                         part, last, 0) == -1)
            return -1;
        offset += part;
    } while (offset < length);
    return 0;
}
\end{lstlisting}
Флаг \code{last} обеспечивает один перевод строки после всех её частей. Разбиение производится по байтам; английские гласные в фильтре представлены однобайтовыми ASCII-символами.

\subsection{Приём и обработка в дочернем процессе}
Ниже показаны основные действия цикла из \code{child.c}. Проверки ошибок вывода сокращены.
\begin{lstlisting}[language=C]
if (wait_for_data(&data->can_read) == -1) {
    result = 1;
    break;
}
if (data->finished) {
    if (sem_post(&data->can_write) == -1) {
        perror("sem_post");
        result = 1;
    }
    break;
}
/* ... */
size_t length = remove_vowels(data->text, data->length);
if (length > 0) {
    if (fwrite(data->text, 1, length, stdout) != length) {
        /* ... */
    }
}
if (data->end_of_line) {
    if (fputc('\n', stdout) == EOF) { /* ... */ }
}
if (sem_post(&data->can_write) == -1) { /* ... */ }
\end{lstlisting}

\subsection{Фильтрация и освобождение ресурсов}
Функция из \code{filter.c} выполняет фильтрацию без выделения второго буфера:
\begin{lstlisting}[language=C]
size_t remove_vowels(char* text, size_t length) {
    size_t j = 0;
    for (size_t i = 0; i < length; i++) {
        if (!is_vowel(text[i])) text[j++] = text[i];
    }
    return j;
}
\end{lstlisting}
\code{is_vowel()} сравнивает символ с \code{a,e,i,o,u,A,E,I,O,U}. Цифры, знаки пунктуации и согласные сохраняются в прежнем порядке.

Перед очисткой родитель ожидает завершения каждого ещё не собранного ребёнка через \code{waitpid()}; ожидание повторяется при \code{EINTR}. Проверяются \code{WIFEXITED(status)} и \code{WEXITSTATUS(status)}. Затем вызывается \code{destroy_shared()} для каждой созданной области.

В \code{destroy_shared()} последовательно выполняются следующие вызовы (между ними в исходнике находятся проверки результата):
\begin{lstlisting}[language=C]
sem_destroy(&data->can_write)
sem_destroy(&data->can_read)
munmap(data, sizeof(SharedData))
shm_unlink(name)
\end{lstlisting}
Ребёнок снимает собственное отображение через \code{close_shared()}, где вызывается \code{munmap()}. \code{munmap()} освобождает отображение процесса, а \code{shm_unlink()} удаляет имя объекта. Это разные операции.

\section{Сборка и запуск}
Проект использует C11 и CMake. \code{parent} собирается из \code{parent.c} и \code{shared.c}; \code{child} --- из \code{child.c}, \code{filter.c} и \code{shared.c}. Обе цели связаны с \code{Threads::Threads}; для GNU/Clang включены \code{-Wall -Wextra -Wpedantic}.
\begin{lstlisting}[language=bash]
cmake -S . -B build
cmake --build build
cd build
./parent
\end{lstlisting}
Запуск выполняется из \code{build}, поскольку путь \code{./child} в \code{execvp()} задан относительно текущего каталога.

\newsection{Результаты функционального запуска}
При интерактивном запуске родитель имел PID 52126, дети --- PID 52159 и 52160. Были указаны файлы \code{out1} и \code{out2}. Оба ребёнка вывели сообщения \code{started} и \code{finished}.

Результаты обработки входных строк приведены в таблице. Обозначение $\varnothing$ означает пустую строку.

\begin{center}
\smalltable
\begin{tabularx}{\linewidth}{@{}c c >{\raggedright\arraybackslash}X >{\raggedright\arraybackslash}X@{}}
\toprule
№ & Файл & Входная строка & Результат\\
\midrule
1 & \code{out1} & \code{aaaaaaa} & $\varnothing$\\
2 & \code{out2} & $\varnothing$ & $\varnothing$\\
3 & \code{out1} & $\varnothing$ & $\varnothing$\\
4 & \code{out2} & \code{akdlkasldkalk} & \code{kdlksldklk}\\
5 & \code{out1} & \code{paspofk} & \code{pspfk}\\
6 & \code{out2} & \code{slkfslakjf} & \code{slkfslkjf}\\
7 & \code{out1} & Строка со скобками, буквами и цифрами (см. ниже) & См. точный результат ниже\\
8 & \code{out2} & \code{gggg} & \code{gggg}\\
9 & \code{out1} & \code{;amibfws34} & \code{;mbfws34}\\
\bottomrule
\end{tabularx}
\end{center}
Для строки 7 исходное и полученное значения:
\begin{lstlisting}
[]][[weflkcaonijefj98549jos[apvk[pmdbna
[]][[wflkcnjfj98549js[pvk[pmdbn
\end{lstlisting}
Файл \code{out1} содержит 49 байт и 5 переводов строки, \code{out2} --- 27 байт и 4 перевода строки. Девять входных строк распределены между процессами в соотношении 5/4.

Запуск показывает обработку пустых строк, удаление строчных английских гласных и сохранение согласных, цифр и специальных символов.

\newsection{Анализ трассировки strace}
\subsection{Порядок запуска}
Для исследования системных вызовов использована команда:
\begin{lstlisting}[language=bash]
printf 'trace1\ntrace2\nhello\nworld\n' | strace -f \
  -e trace=process,memory,openat,close,ftruncate,futex,unlink,unlinkat \
  -o trace.log ./parent
\end{lstlisting}
В файлах \code{trace1} и \code{trace2} получены соответственно \code{hll\n} (4 байта) и \code{wrld\n} (5 байт). Гласные из строк \code{hello} и \code{world} удалены.

\subsection{Создание и подключение общей памяти}
Ниже приведены системные вызовы, связанные с созданием и подключением общей памяти.
\begin{lstlisting}
54251 openat(AT_FDCWD, "/dev/shm/lab4_54251_1", O_RDWR|O_CREAT|O_EXCL|O_NOFOLLOW|O_CLOEXEC, 0600) = 3
54251 ftruncate(3, 4176)                = 0
54251 mmap(NULL, 4176, PROT_READ|PROT_WRITE, MAP_SHARED, 3, 0) = 0x7171c1138000
54251 openat(AT_FDCWD, "/dev/shm/lab4_54251_2", O_RDWR|O_CREAT|O_EXCL|O_NOFOLLOW|O_CLOEXEC, 0600) = 3
54251 ftruncate(3, 4176)                = 0
54251 mmap(NULL, 4176, PROT_READ|PROT_WRITE, MAP_SHARED, 3, 0) = 0x7171c1136000
54252 openat(AT_FDCWD, "/dev/shm/lab4_54251_1", O_RDWR|O_NOFOLLOW|O_CLOEXEC) = 3
54252 mmap(NULL, 4176, PROT_READ|PROT_WRITE, MAP_SHARED, 3, 0 <unfinished ...>
54252 <... mmap resumed>)               = 0x7757db6ae000
54253 openat(AT_FDCWD, "/dev/shm/lab4_54251_2", O_RDWR|O_NOFOLLOW|O_CLOEXEC) = 3
54253 mmap(NULL, 4176, PROT_READ|PROT_WRITE, MAP_SHARED, 3, 0) = 0x7f27a4763000
\end{lstlisting}
Родитель 54251 создал два объекта, а процессы 54252 и 54253 открыли по одному объекту каждый. Для всех четырёх отображений указан \code{MAP_SHARED}. Разные адреса одного объекта у родителя и ребёнка подтверждают независимость их виртуальных адресных пространств.

\subsection{Синхронизация и штатная очистка}
Выдержка с ожиданием родителя и пробуждением со стороны ребёнка:
\begin{lstlisting}
54251 futex(0x7171c1138000, FUTEX_WAIT_BITSET|FUTEX_CLOCK_REALTIME, 0, {tv_sec=1794436131, tv_nsec=979383363}, FUTEX_BITSET_MATCH_ANY <unfinished ...>
54252 futex(0x7757db6ae000, FUTEX_WAKE, 1 <unfinished ...>
54251 <... futex resumed>)              = 0
54252 <... futex resumed>)              = 1
\end{lstlisting}
Адреса соответствуют началу отображений первого объекта, где в \code{SharedData} расположен \code{can_write}. С учётом исходного кода эти операции согласуются с ожиданием и освобождением семафора. Не каждый вызов POSIX-семафора требует отдельного системного вызова, поэтому число операций обмена по числу строк \code{futex} не определяется.

Завершающая часть трассы содержит следующие выбранные строки:
\begin{lstlisting}
54252 munmap(0x7757db6ae000, 4176 <unfinished ...>
54252 <... munmap resumed>)             = 0
54253 munmap(0x7f27a4763000, 4176)      = 0
54252 +++ exited with 0 +++
54253 +++ exited with 0 +++
54251 munmap(0x7171c1138000, 4176)      = 0
54251 unlink("/dev/shm/lab4_54251_1")   = 0
54251 munmap(0x7171c1136000, 4176)      = 0
54251 unlink("/dev/shm/lab4_54251_2")   = 0
54251 +++ exited with 0 +++
\end{lstlisting}
Также в трассе присутствуют уведомления \code{SIGCHLD} для обоих детей со значением \code{si_status=0}. Дети сняли собственные отображения и завершились с кодом 0. Родитель снял оба отображения, успешно удалил оба имени через \code{unlink()} и завершился с кодом 0.

При штатном завершении освобождены отображения и удалены имена обоих созданных объектов общей памяти.

\newsection{Итоги проверки и выводы}
\begin{center}
\smalltable
\begin{tabularx}{\linewidth}{@{}>{\raggedright\arraybackslash}p{0.39\linewidth} >{\raggedright\arraybackslash}X@{}}
\toprule
Проверяемое свойство & Основание и результат\\
\midrule
Два дочерних процесса & В интерактивном запуске указаны PID 52159 и 52160; в трассе видны процессы 54252 и 54253.\\
Чередование строк и фильтрация & Файлы \code{out1/out2}: 5 и 4 строки, результаты соответствуют вводу.\\
Пустые строки и строка из гласных & Пустые результаты сохранены как отдельные строки.\\
Две области \code{MAP_SHARED} & Трасса содержит создание и подключение объектов по 4176 байт.\\
Завершение и очистка в трассируемом запуске & Три кода завершения 0; успешные \code{munmap()} и два \code{unlink()}.\\
\bottomrule
\end{tabularx}
\end{center}

В работе реализован обмен строками между родителем и двумя дочерними процессами через именованные объекты общей памяти. Семафоры \code{can_write} и \code{can_read} обеспечивают последовательную передачу порций в каждом канале. Разделение признаков конца строки и конца обмена позволяет сохранять пустые строки и передавать строки частями.

Результаты запуска подтверждают фильтрацию строк и работу двух каналов. Трассировка показывает создание общей памяти, совместные отображения, операции ожидания и пробуждения, а также штатное освобождение ресурсов.
\end{document}
